Studies and trackers of public-sector AI purchasing usually identify AI contracts with English keyword lists, but the accuracy of this measurement in multilingual procurement data, and its effect on downstream analyses, has rarely been quantified. We compared six scalable methods for identifying AI-related notices among 54,160 IT and software contract award notices published in Tenders Electronic Daily (TED) in 2024--2025, of which 6.1\% were in English. A stratified probability sample of 1,199 notices was annotated with a fixed codebook by two large language models, and design-based supervised learning (DSL) was used to estimate prevalence and the association between AI-relatedness and single bidding. Hypotheses and the sample design were pre-specified. A multilingual dictionary achieved the highest F1 (0.63, 95\% CI 0.50--0.78), doubling the recall of the English dictionary (0.56 vs.\ 0.28). Contrary to the plan, design-weighted embedding and TF-IDF classifiers performed worse (F1 0.06 and 0.10). The DSL prevalence of AI-related notices was 2.32\% (CI 1.56--3.21\%), about three times the 0.83\% flagged by the English dictionary. A naive multilingual-dictionary regression indicated that AI-related contracts were 4.9 percentage points less likely to receive a single tender, but design-based estimates were too imprecise to confirm any association, largely because a few contested labels among unflagged notices carried large weights. The results indicate that keyword-based AI indicators in procurement require multilingual terms and quantified measurement error before they are used for inference.
